\subsection{Преломление}

%%%%%%%%%%%%%%%%%%%%%%%%%%%%%%%%%%%%%%%%%%%%%%%%%%%%
%%%%%%%%%%%%%%%%%%%%%%%%%%%%%%%%%%%%%%%%%%%%%%%%%%%%
 \z{Луч света падает под углом $\alpha$ (к нормали) на стеклянную пластинку толщины $h$ с показателем преломления $n$. 
Докажите, что вышедший луч будет параллелен падающему и найдите расстояние между двумя прямыми, вдоль которых идут эти лучи.}
%
{$s = \frac{h \sin(\alpha - \beta)}{\cos\beta}, \quad \text{где} \quad \sin\beta = \frac{\sin\alpha}{n}$}

%%%%%%%%%%%%%%%%%%%%%%%%%%%%%%%%%%%%%%%%%%%%%%%%%%%%
 \z{На дне сосуда, наполненного водой до высоты $H$, находится точечный источник света $S$. На поверхности воды плавает диск так, что его центр $O$ находится над точечным источником света.  При каком минимальном радиусе диска ни один луч не выйдет через поверхность воды? Показатель преломления воды $n$.}
%
{$R = \frac{h}{\sqrt{n^2 - 1}}$}

%%%%%%%%%%%%%%%%%%%%%%%%%%%%%%%%%%%%%%%%%%%%%%%%%%%%
\z{Точечный источник света расположен на дне водоёма глубиной $h = 0.6\,\text{м}$. В некоторой точке поверхности воды вышедший в воздух преломлённый луч оказался перпендикулярным лучу, отражённому обратно в воду. На каком расстоянии $L$ от источника на дне водоёма достигнет дна отражённый луч? Показатель преломления воды $n = \frac{4}{3}$.}
%
{$L = \frac{2h}{n} = 0.9\,\text{м}$}

%%%%%%%%%%%%%%%%%%%%%%%%%%%%%%%%%%%%%%%%%%%%%%%%%%%%
 \z{Луч света падает на плоскопараллельную стеклянную пластину толщины $d = 2\,\text{см}$ под углом $\alpha = 30\degree$. Какое расстояние $a$ будет между лучами: прошедшим пластину без отражения и претерпевшим двукратное отражение от её граней? Показатель стекла $n = 1.5$.}
%
{$a = \frac{d \,(2\sin\alpha \cos\alpha)}{\cos\beta} = 1.22\,\text{см}, \quad \text{где} \quad \sin\beta = \frac{\sin\alpha}{n}$}

%%%%%%%%%%%%%%%%%%%%%%%%%%%%%%%%%%%%%%%%%%%%%%%%%%%%
 \z[{\begin{wrapfigure}[5]{r}{0.15\textwidth}%
		\vspace{-15pt}%
		\begin{flushright}%
 			\includegraphics[width=0.15\textwidth]{4.3.5}%
		\end{flushright}%
	\end{wrapfigure}%
	}]%
{На равнобедренную стеклянную призму падает широкий параллельный пучок света, перпендикулярный грани $BC$, ширина которой $d = 5\,\text{см}$. 
На каком расстоянии $L$ от грани $BC$ преломлённый призмой свет разделится на два неперекрывающихся пучка? 
Показатель преломления стекла $n = 1.5$, угол при основании призмы $\alpha = 5.7\degree$. 

\noindent
При расчётах учесть, что для малых углов $\tan\alpha \approx \sin\alpha \approx \alpha$.}
%
{$L = \frac{d}{2\alpha (n - 1)}$}

%%%%%%%%%%%%%%%%%%%%%%%%%%%%%%%%%%%%%%%%%%%%%%%%%%%%
 \z{Световой луч падает на поверхность стеклянного шара. Угол падения равен $\alpha = 45\degree$, показатель преломления стекла $n = 1.41$. Найти угол $\gamma$ между падающим лучом и лучом, вышедшим из шара.}
%
{$\gamma = 2\alpha - 2\beta \approx 30\degree, \quad \text{где} \quad \sin\beta = \frac{\sin\alpha}{n}$}

%%%%%%%%%%%%%%%%%%%%%%%%%%%%%%%%%%%%%%%%%%%%%%%%%%%%
 \z{Человек смотрит на рыбку, находящуюся в диаметрально противоположной от него точке шарового аквариума радиусом $R$. На сколько смещено при этом изображение рыбки относительно самой рыбки? Показатель преломления воды $n = \frac{4}{3}$.}
%
{$R$}


